\newcommand{\labname}{18-500 Design Review Report Template - 18 January 2022}
\RequirePackage{rotating} % to get the sidewaysfigure command for rotating figures 90 degrees
\documentclass[10pt]{article}
\usepackage[left=0.65in, right=0.65in, top=0.7in, bottom=0.7in]{geometry}
\usepackage{amsmath,amsfonts,amssymb}
\usepackage{graphicx}
\usepackage{fancyhdr}
\usepackage{tikz}
\usepackage{multicol}
%\usepackage{pdflscape} 
\usepackage{float} % To get the H option on figures
\usepackage[backend=bibtex,
style=numeric,
bibencoding=ascii
]{biblatex}
\addbibresource{sample.bib}
\graphicspath{ {./img/} }
\usepackage{indentfirst}

% LSTLISTINGS PACKAGE (for pseudocode) ======================================
\usepackage{listings}
% END LSTLISTINGS PACKAGE ===================================================

% header
\pagestyle{fancy}
\fancyhf{}
\lhead{\labname}
\rhead{Page \thepage{} of \pageref{mylastpagelabel}}

% tables and figures in column
\makeatletter
\newenvironment{tablehere}
  {\def\@captype{table}}
  {}

\newenvironment{figurehere}
  {\def\@captype{figure}}
  {}
\makeatother
%---------------------------------------------------------------------------------------
%---------------------------------------------------------------------------------------
\begin{document}
\begin{center}
    \Huge{\textbf{Capstone Project Name}} \\
    \vspace*{0.15in}
    \normalsize{Authors: FirstName LastName, FirstName LastName, FirstName LastName} \\
    \normalsize{Affiliation: Electrical and Computer Engineering, Carnegie Mellon University}
\end{center}
\vspace*{0.05in}
\begin{multicols*}{2}
{\small 
    \textbf{
        \emph{Abstract}---A system capable of $\ldots$ See both the guidelines on Canvas as well as the Word template for additional information. \\
        } 
\par}

{\small
    \textbf{
        \emph{Index Terms}---Design, Robot, Sensor (key words or phrases in alphabetical order, use your own words, delete the ones listed as an example) \\
    }
\par}


%---------------------------------------------------------------------------------------
%---------------------------------------------------------------------------------------
\section{INTRODUCTION} \label{introduction}
See both the guidelines on Canvas as well as the Word template for additional information.

You must write something for each of the sections in this template. The Word template has some guidelines on the recommended length for each section.

All teams must either use this Latex template or the Word template provided.

pdflatex template.tex on linux.andrew.cmu.edu compiles this template into a PDF file. Don't forget to run bibtex as it suggests $\ldots$

{\color{red} Differences between the Design Review and Final Report templates are presented using this red highlight.}

%---------------------------------------------------------------------------------------
%---------------------------------------------------------------------------------------
\section{USE-CASE REQUIREMENTS} \label{use-case-requirements}
Based on the application or problem you have described in the Introduction (section \ref{introduction}) $\ldots$


%---------------------------------------------------------------------------------------
%---------------------------------------------------------------------------------------
\section{ARCHITECTURE AND/OR PRINCIPLE OF OPERATION} \label{architecture}
Describe your system architecture and/or principle of operation, but do not be redundant with the System Implementation section \ref{system-implementation}. You will probably need to use figures in this section. 

The next few lines in the Latex template provides a simple example of a two image figure.

\begin{figure}[H]
\vspace{0.5 in}
%\includegraphics{img/image1.png}

{\par\centering (a) \par}
\vspace{0.5 in}

%\includegraphics{img/image2.png}
{\par\centering (b) \par}
\caption{System description. (a) overall system. (b) zoom-in of xxxx.}
\end{figure}


The next few lines in the Latex template mimics the figures in the Word template. Please view the word template for a description about block diagrams and the importance of making sure your figures are readable.
\begin{figure}[H]
\includegraphics{img/Fig1_NotReadable.png}
\caption{An unreadable figure.  Consider making it wide enough to span both columns.}
\label{fig:unreadable-block-diagram} 
\end{figure}

\begin{figure*}
\centering
\includegraphics[width=0.75\textwidth]{img/Fig3_Readable.png}
\caption{This same figure is placed spanning both columns (actually spanning 0.75 of the two columns).}
\label{fig:readable-block-diagram} 
\end{figure*}

If neither Fig.~\ref{fig:unreadable-block-diagram} nor Fig.~\ref{fig:readable-block-diagram} are readable, you can always use Fig.~\ref{fig:full-size-block-diagram} at the end of the document.

%---------------------------------------------------------------------------------------
%---------------------------------------------------------------------------------------
\section{DESIGN REQUIREMENTS} \label{design-requirements}
Please see the Word template and the guidelines on Canvas for more details about mapping the Use-Case Requirements in section \ref{use-case-requirements} to this section on Design Requirements.

%---------------------------------------------------------------------------------------
%---------------------------------------------------------------------------------------
\section{DESIGN TRADE STUDIES} \label{design-trade-studies}
Please see the Word template and the guidelines on Canvas for more details about trade studies.

Trade studies of sub-systems can also be included in this section. You should use 
sections with the \texttt{\\subsection} command to split up this section as dictated
by content.

%---------------------------------------------------------------------------------------
\subsection{Design Specification or Subsystem}
As shown in (1)$\ldots$  which matches the code described in \cite{NACE} we $\ldots$ 


\subsection{Another Design Specification or Subsystem}
Equation (2) $\ldots$, Continue with the next design specification B. Or continue with the trade study for subsystem B.

\subsection{Equation Formatting} \label{equation-formatting}
The equations are an exception to the prescribed specifications of this template. 
Getting equations to look good is one of the major reasons to use \LaTeX~in the first
place.  Use the \texttt{equation} environment so that your equations are
sequentially numbered.
\begin{equation}
    \alpha + \beta = \gamma \label{eq:alpha}
\end{equation}
where $\alpha$ is the first variable, $\beta$ is the second variable and $\gamma$ is the result. All variables are in mks units.

Be sure that the symbols in your equation have been defined before or immediately following the equation. Use ``(\ref{eq:alpha})”, not ``Eq.~(\ref{eq:alpha})” or ``equation~(\ref{eq:alpha})”, except at the beginning of a sentence: ``Equation~(\ref{eq:alpha}) is . . .”

%---------------------------------------------------------------------------------------
\subsection{Figure and Table Formatting}
Place figures and tables at the top and bottom of columns. Avoid placing them in the 
middle of columns. Large figures and tables may span across both columns. Figure 
captions should be below the figures; table heads should appear above the tables. 
Always locate figures and tables after they are cited in the text. Use the abbreviation 
“Fig. 1”, even at the beginning of a sentence.

Figure Labels: Use words rather than symbols or abbreviations when writing Figure 
axis labels to avoid confusing the reader. As an example, write the quantity
“Magnetization”, or “Magnetization, M”, not just “M”. If including units in the 
label, present them within parentheses. Do not label axes only with units. 
In the example, write “Magnetization (A/m)” not just “A/m”. Do not label axes 
with a ratio of quantities and units. For example, write “Temperature (K)”, 
not “Temperature/K”.

%---------------------------------------------------------------------------------------
%---------------------------------------------------------------------------------------
\section{SYSTEM IMPLEMENTATION} \label{system-implementation}
Again please use the guidelines on Canvas and the Word template for what to include in this section. This section should be complementary in content with the Architecture 
section \ref{architecture} rather than redundant. You can refer back to earlier figures in section \ref{architecture} using Fig.~\ref{fig:unreadable-block-diagram} and Fig.~\ref{fig:readable-block-diagram}.

There should be a subsection for each of the subsystems as shown below.
\subsection{Subsystem A}
Description of subsystem A (often helps to include a figure).

\subsection{Subsystem B}
Description of subsystem B (often helps to include a figure).

%---------------------------------------------------------------------------------------
%---------------------------------------------------------------------------------------
\section{TEST \& VALIDATION} \label{test-validation}
Again please use the guidance on Canvas and the Word template for what to include in this section.

\subsection{{\color{red}Results} for Design Specification A}
Planned test {\color{red} has been replaced by result}.

\subsection{{\color{red}Results} for Design Specification B}
Continue with testing and validation results for the next design specification.

%---------------------------------------------------------------------------------------
%---------------------------------------------------------------------------------------
\section{PROJECT MANAGEMENT} \label{project-management}
\subsection{Schedule} \label{schedule}
Again please use the guidance on Canvas and the Word template for what to include here.
The schedule is shown in Fig.~\ref{fig:gantt-chart}.

%---------------------------------------------------------------------------------------
\subsection{Team Member Responsibilities}
Again please use the guidance on Canvas and the Word template for what to include here.

%---------------------------------------------------------------------------------------
\subsection{Bill of Materials and Budget} \label{budget}
Again please use the guidance on Canvas and the Word template for what to include here.

Here is a simple example of a Latex table which you can include in the extra page provided for the BOM and you can refer to as Table~\ref{budget-table}.
\begin{table*} 
  \centering
  \caption{Bill of materials} \label{budget-table}
  \vspace{6 pt}
  \begin{tabular}{ l c c c r r }
     \hline
     Description & Model \# & Manufacturer & Quantity & Cost @ & Total \\ 
     \hline
     part description 1 & 0001 & MyManf & 2 & \$1.99 & \$3.98 \\  
     part description 2 & 0022 & YourManf & 1 & \$10.02 & \$10.02 \\  
     \hline
     &  &  &  &  & \$14.00 \\  
  \end{tabular}
\end{table*}


%---------------------------------------------------------------------------------------
\subsection{{\color{red}TechSpark Usage}}
{\color{red} See guidance on Canvas and Word template for what to include here.}

%---------------------------------------------------------------------------------------
\subsection{Risk {\color{red}Management}}
Again please use the guidance on Canvas and the Word template for what to include here. {\color{red}Note Mitigation Plans has now become Management}.

%---------------------------------------------------------------------------------------
%---------------------------------------------------------------------------------------
\section{{\color{red}ETHICAL ISSUES}} \label{ethical-issues}
{\color{red}See guidance on Canvas and Word template for what to include in this Final Report specific section}.

%---------------------------------------------------------------------------------------
%---------------------------------------------------------------------------------------
\section{RELATED WORK} \label{related-work}
Again please use the guidance on Canvas and the Word template for what to include here.

%---------------------------------------------------------------------------------------
%---------------------------------------------------------------------------------------
\section{SUMMARY} \label{summary}
Again please use the guidance on Canvas and the Word template for what to include here.
{\color{red}
Also include two new subsections
\subsection{Future work}
See Guidance and Word template.

\subsection{Lessons Learned}
See Guidance and Word template.
}

\section*{Glossary of Acronyms}
Include an alphabetized list of acronyms if you have lots of these included in your document. Otherwise define the acronyms inline.
\begin{itemize}
    \item MQTT – Message Queuing Telemetry Transport
    \item OBD – On-Board Diagnostics
    \item RPi – Raspberry Pi 
\end{itemize}

%---------------------------------------------------------------------------------------
\printbibliography
Your references should be a very carefully crafted list, cited in the appropriate ways.  Don't merely list a Wikipedia page or a bunch of GitHub URLs.  Note that any code you used in your project does need to be cited.
\\

You can insert blank pages after the references to add full page figures or tables for
\begin{itemize}
\item Architecture and system description figures – if so, make sure you refer to them in section \ref{architecture} or \ref{system-implementation} as appropriate.
\item Milestone and Schedule chart – if so, make sure you refer to it from section \ref{schedule}.
\item Budget and Parts list – if so make sure you refer to it from section \ref{budget}.
\end{itemize}

You are allowed no more than 3 optional pages to ensure any large system diagrams, and your milestones and budget are readable at the end of your document.


% This label is used in the header to keep track of the page count.
\label{mylastpagelabel}

\end{multicols*}

\begin{sidewaysfigure}
  \includegraphics[height=7.1 in,width=\textwidth,keepaspectratio]{img/Fig3_Readable.png}
  \caption{A full-page version of the same system block diagram as depicted earlier.}
  \label{fig:full-size-block-diagram} 
\end{sidewaysfigure}

\begin{sidewaysfigure}
  \includegraphics[height=7.1in,width=\textwidth,keepaspectratio]{img/Gantt.png}  
  \caption{Gantt Chart}
  \label{fig:gantt-chart} 
\end{sidewaysfigure}

\end{document}





